% What happens when the Guiding Principles cease to be self-guiding?

% Note that Anderson is considered required reading should the AFLC every think to change the principles.

% This is a fundamental issue: "Principles intended to show us the way in congregational work ought surely to be so readily understandable that we can understand what is meant without having to use a commentary."


\section{Answer}

(L. O. Anderson.)

“Folkebladet,” August 7, 1918


——

When, in Folkebladet for July 17, I am requested by Prof. Harbo to give a more precise explanation of wherein consists the difference between the view of congregational work which the Lutheran Free Church maintains through the Leading Principles and the view of congregational work which the large church bodies maintain, I will say, first of all, that it ought not to be necessary for me to have to give Prof. Harbo a more detailed explanation of a matter which, more than almost any other matter, has been discussed and explained again and again by men within the Lutheran Free Church, in speech and in writing, ever since the Lutheran Free Church began its existence as an independent church body. It would be altogether wasted time for me to attempt it. In my previous article I pointed out the difference, and I will not enter further into it on the present occasion.

That there is a difference—a considerable difference—between the view of congregational work maintained in the Leading Principles and that maintained in the church bodies’ rules for congregational work is something that Free Church people, least of all, will dispute. If there is no difference—and that an essential difference—then the Free Church congregations have no valid excuse for having refrained from joining the union last year.

---

Upon looking more closely, one will find that in my previous article I do not make the expression, the view of congregational work maintained in the Leading Principles, identical with the expression, the view of congregational work maintained in the Lutheran Free Church, as my esteemed teacher does in his article. To make these two expressions identical would not come sufficiently near the truth. For a large part of those who profess adherence to the Free Church—I among them—do not carry into effect in congregational work what is maintained on this point in the Leading Principles. We both organize congregations out of worldly people and admit into the congregation persons who do not profess themselves to be believers. Is this right according to the Leading Principles?

But, it is said, we must have an external organization if we are to make congregational work function, and this is in harmony with the Leading Principles, point 3. Yes, but point 2 emphasizes that “the congregation consists of believing persons who, through the use of the means of grace and the gifts of grace according to the direction of God’s Word, seek their own salvation and eternal blessedness and that of all their fellow human beings.” Someone will immediately remark that by this is meant the congregation in the spiritual sense, the true believers, whom alone the Lord acknowledges as belonging to the congregation of God. Then go to point 4. And let us take it as it stands. Principles intended to show us the way in congregational work ought surely to be so readily understandable that we can understand what is meant without having to use a commentary. In point 4 it says:

“Since not all members of the congregation’s external organization are always believing persons, and since such hypocrites often take false comfort from their connection with the congregation, it is the congregation’s holy calling to cleanse itself both through awakening preaching of God’s Word, through earnest admonition and reminder, and through the exclusion of open and obstinate sinners.”

Does it follow from point 4 that we have the right to establish a congregation out of worldly people? Does it appear to be in harmony with the Leading Principles that we deliberately admit non-believing persons—that is, “hypocrites”—into the congregation? Have not the church bodies been referred to as those who conduct congregational work in this manner? And has there not, especially in the Free Church’s earliest period, been earnest warning against this, which has been called the “mass congregation”?

What, then, are we to do? Do we act rightly toward the Free Church and toward ourselves by continuing to remain connected with the Free Church, while at the same time continuing deliberately to transgress the Leading Principles in so important a field as the work of the local congregation? To approve this—would it not be like wrenching out one of the foundation stones of a building?—an act that would render the whole building insecure. Is it not time that we who do not carry the Leading Principles into practice either turn back and begin to practice what we have professed—so that we may no longer have to expose ourselves to the charge of sailing under false colors—or else leave the Free Church and go where, according to our practice, we belong?

---

I have chosen the latter. In saying this I also say, of course, that I do not believe in carrying the Leading Principles through in congregational work. From what has been stated above, it is not difficult to guess the reason.

In the Leading Principles, point 4, those persons who enter the congregation’s external organization and are not believers are branded as hypocrites. (Would not point 4 have been much better in the form in which it originally stood, as it came from Sverdrup’s pen? See Guidance in the Principles of the Lutheran Free Church, page 50.) I say nothing against the latter part of point 4, that it is “the congregation’s holy calling to cleanse itself both through awakening preaching of God’s Word, through earnest admonition and reminder, and through the exclusion of open and obstinate sinners.” That is biblical. But is it also biblical to brand persons as hypocrites when they join the congregation’s external organization and are not believing persons? It is, according to Pontoppidan. But is it biblical?

I have found nothing in the Bible that urges me to brand as a hypocrite the person who, without being a believer, requests admission into the congregation’s external organization. But I have found passages in the Bible that encourage me rather to regard the desire of non-believing persons to join the congregation’s external organization as an expression of an inward longing for the truth. And precisely the longing for the truth is the point of contact between the believing and the non-believing soul, and ought to make us cautious, so that neither in speech nor in writing, neither in doctrinal statement nor in principles for congregational work, do we give it to be understood that such a person is a hypocrite. The Savior never gave any such thing to be understood when he dealt with non-believing persons. Nor did the apostles! It was persons who outright hardened themselves against the truth and who nevertheless wished to present themselves before others as pious whom the Lord and the apostles called hypocrites. (That hypocrites are found in the congregation among those who profess themselves to be believers is granted, and shall not be discussed further here.)

My experience in congregational work confirms what has been stated above. Where there is “awakening preaching of God’s Word, earnest admonition and reminder,” there “obstinate sinners” will not enter the congregation. It is impossible to get them to join, whether one attempts to entice them or to use force. They will sooner let the veil slip aside and show themselves in their true form. Furthermore, several who gave the impression of being thoroughly worldly when they enrolled in the congregation have since, through the working of the Word and the Holy Spirit, become earnest believers, zealous for the Lord’s cause.

---

But have you laid the Free Church’s principles before your congregation, and have you attempted to train the people so that they might be able to acknowledge the principles and carry them into practice?

The fact is that our congregation is one of the many congregations that have professed adherence to the Free Church but have not, by congregational resolution, acknowledged the Free Church’s principles and rules and reported this to the secretary of the organization committee (see Revised Rules, point 4), and which consequently, strictly speaking, have not belonged to the Free Church. But how could I have done so? It would have been meaningless to attempt to train the people of the congregation to acknowledge and carry into practice something directly contrary to the lines along which the congregation had been organized. The congregation had, after all, been established before I came, without the slightest regard for the Leading Principles. They were not even read aloud. Had we attempted to secure acknowledgment of the principles, it would certainly have caused an upheaval, to the injury of the congregational work here.

I regret that it took me a full three years to come to a clear understanding of the hollowness of this procedure and to learn to act consistently. We must be consistent in congregational work if we are to expect the Lord to bless it.

The translation retains “mass congregation,” “awakening preaching,” “external organization,” and “Leading Principles” as historically and theologically marked terminology rather than replacing them with smoother modern ecclesiastical equivalents.
